\documentclass[b5paper,11pt,twoside,icelandic]{article}
\input{preamble}

\title{Titill}
\author{Höfundur}

\begin{document}
\maketitle

\begin{abstract}
  Ágrip hér
\end{abstract}

\section{Kafli}

Steven Krantz hefur skrifað bók \cite{Kra:1997}
um það hvernig á að skrifa
stærðfræðitexta sem er einnig hægt að nálgast á netinu, 
\url{https://arxiv.org/pdf/1612.04888v1.pdf}. 
Martin Erickson hefur einnig skrifað stutta grein sem getur gagnast
ykkur,
\url{http://www.xavier.edu/math-undergraduate-research/documents/Write.pdf}.
Hinn þekkti stærðfræðingur Paul Halmos skrifaði grein um sama efni, \url{http://www.math.washington.edu/~lind/Resources/Halmos.pdf}, en
Krantz tileinkar bókina sína einmitt Halmos.



\subsection{Undirkafli} 

Texti.

\begin{definition}\label{skilgr}
  Skilgreinum \emph{blah} sem bleh.
\end{definition}

\begin{lemma}
Fyrir $x \in \mathcal E$ þá er
\begin{equation}\label{jafna}
 \overline x \to S_n
 \end{equation}
 \end{lemma}
	

\begin{theorem}
  Setning
\end{theorem}

\begin{proof}
 \todo{ToDo: Athugasemd}
Sönnun. Ef blah (sjá \ref{skilgr}) þá er
  \[
     S_n = \frac{X_1 + X_2 + \cdots + X_n}{n}
         = \frac{1}{n}\sum_{i}^{n} X_i
  \]
  samkvæmt \eqref{jafna}

\end{proof}

\begin{center}
  \includegraphics[width=4cm]{./frog.jpg}
\end{center}

Listi yfir heimildir er geymdur í skránni \verb+bibtex.bib+. 
Auðveldasta leiðin til að bæta í hana er að finna heimildina
á \url{http://ams.org/mathscinet}. Smella á \emph{Select 
alternative format}, velja \emph{Bibtex} og líma textann sem
birtist í skrána.

Frekari leiðbeining um \LaTeX má finna í 
\emph{The Not So Short Introduction to LaTeX}, 
\url{https://tobi.oetiker.ch/lshort/lshort.pdf},
og á \url{https://en.wikibooks.org/wiki/LaTeX}.

\bibliographystyle{siam}
\bibliography{bibtex}

\end{document}
